% Generated draft: requires mathematical and source review.
\documentclass[11pt]{article}
\usepackage[margin=25mm]{geometry}
\usepackage{fontspec}
\setmainfont{Latin Modern Roman}
\usepackage{amsmath,amssymb,mathtools,unicode-math}
\setmathfont{Latin Modern Math}
\usepackage{enumitem,needspace,xurl,hyperref}
\setlength{\parindent}{0pt}
\setlength{\parskip}{5pt}
\setlength{\emergencystretch}{3em}
\newcommand{\N}{\mathbb{N}}
\newcommand{\Z}{\mathbb{Z}}
\newcommand{\Q}{\mathbb{Q}}
\newcommand{\R}{\mathbb{R}}
\newcommand{\C}{\mathbb{C}}
\newcommand{\F}{\mathbb{F}}
\newcommand{\A}{\mathbb{A}}
\newcommand{\PP}{\mathbb{P}}
\newcommand{\E}{\mathbb{E}}
\newcommand{\eps}{\varepsilon}
\newcommand{\dd}{\,\mathrm{d}}
\begin{document}
\section*{UnsolvedMath ID: 30000480}
\textbf{Source problem code:} catalogue entry
\subsection{Definitions and mathematical statement}
Let $A$ be a Hopf algebra over the field $k$ intended by the source. Thus $A$ has multiplication $m$, unit $\eta$, coproduct $\Delta$, counit $\epsilon$, and antipode $S:A\to A$ satisfying, in particular,
\[
 m(S\otimes\operatorname{id})\Delta
   =\eta\epsilon
   =m(\operatorname{id}\otimes S)\Delta.
\]
Bijectivity of $S$ is not included in the definition of Hopf algebra used for this question.

Assume $A$ is noetherian as an algebra. Under the usual two-sided convention, every ascending chain of left ideals and every ascending chain of right ideals stabilizes. The question is whether
\[
 \forall A\text{ noetherian Hopf algebra}:\qquad
 S:A\longrightarrow A\text{ is bijective}.
\]
Bijective means both injective and surjective as a map, equivalently as a $k$-linear map; it does not mean that the antipode is an algebra homomorphism rather than an antihomomorphism.

The source does not specify the base field, its characteristic, or whether noetherian is intended on one or both sides. The draft states the standard two-sided reading while retaining that convention issue. No finite-dimensionality, affine finite generation, semiprimeness, polynomial-identity assumption, or homological regularity is added. The supplied assessment mentions results under additional hypotheses, but those hypotheses are contextual and do not narrow the universal question.
\subsection{Short English statement}
Must the antipode of every noetherian Hopf algebra be both injective and surjective? Bijectivity is part of the proposed conclusion, with the noetherian side convention and base field still needing source confirmation.
\subsection{Sources}
\begin{itemize}
\item Local UnsolvedMath registry, permanent ID 30000480, code catalogue entry, ``Bijectivity of Antipodes in Noetherian Hopf Algebras.'' Proposed year 2006; proposer unspecified.
\item \emph{Interactions between Algebraic Geometry and Noncommutative Algebra} (2007), Oberwolfach Reports, pp.\ 1319--1384, DOI 10.4171/owr/2006/23; original Report 23/2006 \url{https://doi.org/10.4171/owr/2006/23}. The displayed report year 2007 differs from the DOI and proposed year 2006.
\item Supplied reference ``Lv--Oh--Wang--Yu, bijectivity of the antipode'': \url{https://arxiv.org/abs/1704.05731}.
\item Statement status corrected\_verified; the supplied note says official-report content was checked and notation, context, and question form normalized without changing the source claim.
\item Literature triage checked 2026-08-21: open. It reports positive results requiring extra conditions such as semiprimeness, PI, or homological assumptions, but does not state their exact hypotheses. These are dated supplied claims requiring independent verification. Created 2026-08-25; updated 2026-09-16.
\item Registry curation by Ulam AI / UnsolvedMath Contributors, CC BY 4.0; underlying report and literature provenance remain as supplied. Unreviewed statement-only draft.
\end{itemize}
\end{document}
